% GENERATED FILE. Edit canonical lessons, then run npm run generate:books.
% canonical-source-hash: fnv1a64:b6ede5bd9451453d
% canonical-lessons: GU-C252-pharaj, GU-C252-javabdari, GU-C252-acarya, GU-C252-ganvesh, GU-C252-nivrtt-thavu

\chapter{Duties and responsibility}
\label{ch:gu-duties-and-responsibility}
\hlchaptermodality{252}

\begin{chapteropening}
\textbf{By the end of this chapter:} \emph{I can talk about duty, responsibility and retiring from work.}

Complete the last lesson of chapter 252: I can talk about duty, responsibility and retiring from work.
\end{chapteropening}

\section[\texorpdfstring{\gu{ફરજ} (pharaj)}{ફરજ (pharaj)}]{\gu{ફરજ} (pharaj) — duty}

\label{lesson:GU-C252-pharaj}



\begin{quote}
Before the new one: say the Gujarati for to pay attention, then the Gujarati for to memorize.
\end{quote}

\subsection*{You'll want to know: \gu{ફરજ}}
\textbf{\gu{ફરજ}} — \emph{pharaj} — \textquotedblleft{}duty\textquotedblright{}.

\begin{quote}
\textbf{\gu{આ} \gu{મારી} \gu{ફરજ} \gu{છે}.} — \emph{ā mārī pharaj che} — \textquotedblleft{}This is my duty.\textquotedblright{}
\end{quote}

\begin{cousinweb}
\textbf{\gu{ફરજ}} (\emph{pharaj}) is a loan from Arabic \emph{farḍ}, \textquotedblleft{}obligation\textquotedblright{}, which reached Gujarati through Persian \emph{farz} (EtymDB). Its Arabic root is \emph{f-r-ḍ} (Lane).
\end{cousinweb}

\begin{grammarlens}[title={What you've built}]
One more word for everyday life at A2.
\end{grammarlens}

\subsection*{Your turn}
\begin{itemize}
\raggedright
  \item \emph{Say it:} \emph{pharaj}
  \item \emph{Say it:} \emph{pharaj}, once more
  \item \emph{Say it:} \emph{ā mārī pharaj che} — This is my duty.
  \item \emph{Recall:} say the Gujarati for to pay attention, then the Gujarati for to memorize, then say \emph{pharaj} again
\end{itemize}

\begin{tcolorbox}[breakable,colback=teal!4,colframe=teal!35!black,title={Before you move on}]
What does \gu{ફરજ} mean? (\textquotedblleft{}Duty\textquotedblright{}.) Say it once more.
\end{tcolorbox}
\section[\texorpdfstring{\gu{જવાબદારી} (javābdārī)}{જવાબદારી (javābdārī)}]{\gu{જવાબદારી} (javābdārī) — responsibility}

\label{lesson:GU-C252-javabdari}



\begin{quote}
Before the new one: say the Gujarati for to memorize, then the Gujarati for duty.
\end{quote}

\subsection*{You'll want to know: \gu{જવાબદારી}}
\textbf{\gu{જવાબદારી}} — \emph{javābdārī} — \textquotedblleft{}responsibility\textquotedblright{}.

\begin{quote}
\textbf{\gu{ઘરની} \gu{જવાબદારી} \gu{મારા} \gu{પિતા} \gu{પાસે} \gu{છે}.} — \emph{gharnī javābdārī mārā pitā pāse che} — \textquotedblleft{}My father has the responsibility for the house.\textquotedblright{}
\end{quote}

\begin{cousinweb}
\textbf{\gu{જવાબદારી}} (\emph{javābdārī}) is built on \emph{javābdār}, \textquotedblleft{}answerable\textquotedblright{}, with \emph{-ī} (EtymDB). \emph{Javāb}, \textquotedblleft{}answer\textquotedblright{}, is a loan from Arabic \emph{jawāb}, root \emph{j-w-b} (Lane), and \emph{-dār}, \textquotedblleft{}holder\textquotedblright{}, is usually traced to Persian.
\end{cousinweb}

\begin{grammarlens}[title={What you've built}]
One more word for everyday life at A2.
\end{grammarlens}

\subsection*{Your turn}
\begin{itemize}
\raggedright
  \item \emph{Say it:} \emph{javābdārī}
  \item \emph{Say it:} \emph{javābdārī}, once more
  \item \emph{Say it:} \emph{gharnī javābdārī mārā pitā pāse che} — My father has the responsibility for the house.
  \item \emph{Recall:} say the Gujarati for to memorize, then the Gujarati for duty, then say \emph{javābdārī} again
\end{itemize}

\begin{tcolorbox}[breakable,colback=teal!4,colframe=teal!35!black,title={Before you move on}]
What does \gu{જવાબદારી} mean? (\textquotedblleft{}Responsibility\textquotedblright{}.) Say it once more.
\end{tcolorbox}
\section[\texorpdfstring{\gu{આચાર્ય} (ācārya)}{આચાર્ય (ācārya)}]{\gu{આચાર્ય} (ācārya) — principal, head teacher}

\label{lesson:GU-C252-acarya}



\begin{quote}
Before the new one: say the Gujarati for duty, then the Gujarati for responsibility.
\end{quote}

\subsection*{You'll want to know: \gu{આચાર્ય}}
\textbf{\gu{આચાર્ય}} — \emph{ācārya} — \textquotedblleft{}principal, head teacher\textquotedblright{}.

\begin{quote}
\textbf{\gu{આચાર્ય} \gu{દરેક} \gu{વર્ગમાં} \gu{જાય} \gu{છે}.} — \emph{ācārya darek vargmā̃ jāy che} — \textquotedblleft{}The principal goes to every class.\textquotedblright{}
\end{quote}

A masculine noun.

\begin{cousinweb}
\textbf{\gu{આચાર્ય}} (\emph{ācārya}) is Sanskrit \emph{ācārya}, \textquotedblleft{}teacher\textquotedblright{}, taken over whole as a tatsama. CDIAL has no entry for it, though it lists the related \emph{ācāra}, \textquotedblleft{}mode of action, conduct\textquotedblright{} (CDIAL 1071).
\end{cousinweb}

\begin{grammarlens}[title={What you've built}]
One more word for everyday life at A2.
\end{grammarlens}

\subsection*{Your turn}
\begin{itemize}
\raggedright
  \item \emph{Say it:} \emph{ācārya}
  \item \emph{Say it:} \emph{ācārya}, once more
  \item \emph{Say it:} \emph{ācārya darek vargmā̃ jāy che} — The principal goes to every class.
  \item \emph{Recall:} say the Gujarati for duty, then the Gujarati for responsibility, then say \emph{ācārya} again
\end{itemize}

\begin{tcolorbox}[breakable,colback=teal!4,colframe=teal!35!black,title={Before you move on}]
What does \gu{આચાર્ય} mean? (\textquotedblleft{}Principal, head teacher\textquotedblright{}.) Say it once more.
\end{tcolorbox}
\section[\texorpdfstring{\gu{ગણવેશ} (gaṇveś)}{ગણવેશ (gaṇveś)}]{\gu{ગણવેશ} (gaṇveś) — uniform}

\label{lesson:GU-C252-ganvesh}



\begin{quote}
Before the new one: say the Gujarati for responsibility, then the Gujarati for principal.
\end{quote}

\subsection*{You'll want to know: \gu{ગણવેશ}}
\textbf{\gu{ગણવેશ}} — \emph{gaṇveś} — \textquotedblleft{}uniform\textquotedblright{}.

\begin{quote}
\textbf{\gu{શાળામાં} \gu{બધા} \gu{ગણવેશ} \gu{પહેરે} \gu{છે}.} — \emph{śāḷāmā̃ badhā gaṇveś pahere che} — \textquotedblleft{}Everyone wears a uniform at school.\textquotedblright{}
\end{quote}

A masculine noun.

\begin{cousinweb}
\textbf{\gu{ગણવેશ}} (\emph{gaṇveś}), the dress of a group, is a modern learned compound built from Sanskrit \emph{gaṇa}, \textquotedblleft{}group\textquotedblright{}, and \emph{vēṣa}, \textquotedblleft{}dress\textquotedblright{}; CDIAL lists Old Gujarati \emph{vesa}, \textquotedblleft{}dress\textquotedblright{} (CDIAL 12129).
\end{cousinweb}

\begin{grammarlens}[title={What you've built}]
One more word for everyday life at A2.
\end{grammarlens}

\subsection*{Your turn}
\begin{itemize}
\raggedright
  \item \emph{Say it:} \emph{gaṇveś}
  \item \emph{Say it:} \emph{gaṇveś}, once more
  \item \emph{Say it:} \emph{śāḷāmā̃ badhā gaṇveś pahere che} — Everyone wears a uniform at school.
  \item \emph{Recall:} say the Gujarati for responsibility, then the Gujarati for principal, then say \emph{gaṇveś} again
\end{itemize}

\begin{tcolorbox}[breakable,colback=teal!4,colframe=teal!35!black,title={Before you move on}]
What does \gu{ગણવેશ} mean? (\textquotedblleft{}Uniform\textquotedblright{}.) Say it once more.
\end{tcolorbox}
\section[\texorpdfstring{\gu{નિવૃત્ત} \gu{થવું} (nivṛtt thavũ)}{નિવૃત્ત થવું (nivṛtt thavũ)}]{\gu{નિવૃત્ત} \gu{થવું} (nivṛtt thavũ) — to retire}

\label{lesson:GU-C252-nivrtt-thavu}



\begin{quote}
Before the new one: say the Gujarati for principal, then the Gujarati for uniform.
\end{quote}

\subsection*{You'll want to know: \gu{નિવૃત્ત} \gu{થવું}}
\textbf{\gu{નિવૃત્ત} \gu{થવું}} — \emph{nivṛtt thavũ} — \textquotedblleft{}to retire\textquotedblright{}.

\begin{quote}
\textbf{\gu{મારા} \gu{પિતા} \gu{આ} \gu{વરસે} \gu{નિવૃત્ત} \gu{થાય} \gu{છે}.} — \emph{mārā pitā ā varse nivṛtt thāy che} — \textquotedblleft{}My father retires this year.\textquotedblright{}
\end{quote}

\begin{cousinweb}
\textbf{\gu{નિવૃત્ત} \gu{થવું}} (\emph{nivṛtt thavũ}) is built from \emph{nivṛtt}, Sanskrit \emph{nivṛtta}, \textquotedblleft{}turned back, ceased\textquotedblright{}, taken over as a tatsama, and the verb \emph{thavũ}, \textquotedblleft{}to become\textquotedblright{}. CDIAL has no entry for \emph{nivṛtta}.
\end{cousinweb}

\begin{grammarlens}[title={What you've built}]
One more word for everyday life at A2.
\end{grammarlens}

\subsection*{Your turn}
\begin{itemize}
\raggedright
  \item \emph{Say it:} \emph{nivṛtt thavũ}
  \item \emph{Say it:} \emph{nivṛtt thavũ}, once more
  \item \emph{Say it:} \emph{mārā pitā ā varse nivṛtt thāy che} — My father retires this year.
  \item \emph{Recall:} say the Gujarati for principal, then the Gujarati for uniform, then say \emph{nivṛtt thavũ} again
\end{itemize}

\begin{tcolorbox}[breakable,colback=teal!4,colframe=teal!35!black,title={Before you move on}]
What does \gu{નિવૃત્ત} \gu{થવું} mean? (\textquotedblleft{}To retire\textquotedblright{}.) Say it once more.
\end{tcolorbox}
